This chapter presents the architecture of RACS from the sender's and the receiver's side.
In the diagrams, grey blocks are components inherited unchanged from the baseline
(AgentSemSteCom) and green blocks are new in RACS.

\section{Information Model}
\begin{table}[H]
	\centering
	\caption{What each party knows.}
	\label{tab:info}
	\begin{tabularx}{\textwidth}{|c|c|X|}
		\hline
		\textbf{Symbol} & \textbf{Visibility} & \textbf{Meaning} \\ \hline
		$X_s$ & Private & The secret image \\ \hline
		$K_{\mathrm{pub}}$ & Public & A text description of the image, written by a language model \\ \hline
		$X_{\mathrm{feat}}$ & Public & The image's structure: a segmentation map or a pose skeleton \\ \hline
		$\tau$ & Shared secret & An 8-digit token, the only secret sender and receiver share \\ \hline
		$X_{\mathrm{ref}}$ & Derived & A reference image that both ends must reproduce identically \\ \hline
		$X_{\mathrm{stego}}$ & Transmitted & The innocent-looking image that is actually sent \\ \hline
	\end{tabularx}
\end{table}

Every private random value (seeds, noise, mask, selection) is derived from $\tau$ with SHA-256,
so both ends compute exactly the same values on different machines.

\section{Starting Point: the Baseline}
In AgentSemSteCom, the sender describes the secret in words ($K_{\mathrm{pub}}$) and extracts its
structure ($X_{\mathrm{feat}}$); a diffusion model then paints a reference image $X_{\mathrm{ref}}$
from these and a token-seeded noise; EDICT blends the secret toward that reference to create the
stego image, which is sent through the SwinJSCC semantic codec over a noisy channel. The receiver
paints the same reference again and runs EDICT backwards. This costs about 4\,s and 2.6\,GB of
GPU memory per reference at \emph{each} end, and the secret comes back only approximately.

\section{RACS Architecture}
RACS keeps everything that works in the baseline and changes two things: the reference image is
\emph{retrieved} instead of generated, and an exact-recovery layer carries the secret losslessly
(Path B) next to the semantic channel (Path A). Figures~\ref{fig:racs-tx} and~\ref{fig:racs-rx}
show the sender and the receiver. Green blocks are new, grey blocks are the baseline's own
unchanged code, and the dashed red block is the baseline stage that RACS removes.

\begin{figure}[H]
	\centering
	% Generated from racs_sender.drawio by drawio2tikz.py; edit the .drawio file and rerun.
\begin{tikzpicture}
	\hyphenpenalty=10000\exhyphenpenalty=10000
	\draw[dashed, draw=black!70, rounded corners=2pt] (0.720,-3.420) rectangle (6.120,-8.424);
	\draw[dashed, draw=black!70, rounded corners=2pt] (3.600,-12.672) rectangle (12.240,-17.676);
	\node[draw=black!75, fill=black!6, line width=0.6pt, rounded corners=0.095cm, minimum width=4.680cm, minimum height=1.188cm, text width=4.572cm, align=center, inner sep=1pt, font=\diagfont] (3) at (3.420,-4.554) {\textbf{N1-N2: CLIP search} \\ $K_{\mathrm{pub}}$ $\rightarrow$ CLIP text vector, exact search, 1,000 images $\rightarrow$ top 50};
	\node[draw=black!75, fill=black!6, line width=0.6pt, rounded corners=0.095cm, minimum width=4.680cm, minimum height=1.188cm, text width=4.572cm, align=center, inner sep=1pt, font=\diagfont] (4) at (3.420,-6.102) {\textbf{N3-N4: Rerank} \\ 0.75 meaning + 0.25 structure ($X_{\mathrm{feat}}$ edges) $\rightarrow$ top 10};
	\node[draw=black!75, fill=black!6, line width=0.6pt, rounded corners=0.081cm, minimum width=4.680cm, minimum height=1.008cm, text width=4.572cm, align=center, inner sep=1pt, font=\diagfont] (5) at (3.420,-7.560) {\textbf{N5: Token pick} \\ SHA-256($\tau$) mod 10 $\rightarrow$ $X_{\mathrm{ref}}$};
	\node[draw=black!75, fill=black!6, line width=0.6pt, rounded corners=0.087cm, minimum width=3.960cm, minimum height=1.082cm, text width=3.852cm, align=center, inner sep=1pt, font=\diagfont] (8) at (6.300,-14.701) {\textbf{N9-N10: Hide in busy pixels} \\ add camera-like noise to $X_{\mathrm{stego}}$, $\pm$1 changes in the most textured 50 \%};
	\node[draw=black!75, fill=black!6, line width=0.6pt, rounded corners=0.065cm, minimum width=3.960cm, minimum height=0.811cm, text width=3.852cm, align=center, inner sep=1pt, font=\diagfont] (9) at (6.300,-17.000) {\textbf{N11: Stego PNG file} \\ to receiver (Path B)};
	\node[draw=black!75, fill=black!6, line width=0.6pt, rounded corners=0.270cm, minimum width=2.700cm, minimum height=0.540cm, text width=2.592cm, align=center, inner sep=1pt, font=\diagfont] (10) at (4.950,-0.810) {Secret image $X_{\mathrm{s}}$};
	\node[draw=black!75, fill=black!6, line width=0.6pt, rounded corners=0.270cm, minimum width=2.340cm, minimum height=0.540cm, text width=2.232cm, align=center, inner sep=1pt, font=\diagfont] (11) at (10.530,-0.810) {Private token $\tau$};
	\node[draw=black!75, fill=black!6, line width=0.6pt, rounded corners=0.063cm, minimum width=3.060cm, minimum height=0.792cm, text width=2.952cm, align=center, inner sep=1pt, font=\diagfont] (12) at (3.420,-2.196) {\textbf{$K_{\mathrm{pub}}$ generator} \\ BLIP + Qwen LLM};
	\node[draw=black!75, fill=black!6, line width=0.6pt, rounded corners=0.069cm, minimum width=3.780cm, minimum height=0.864cm, text width=3.672cm, align=center, inner sep=1pt, font=\diagfont] (13) at (7.830,-2.232) {\textbf{Structure extractor} \\ OneFormer / OpenPose $\rightarrow$ $X_{\mathrm{feat}}$};
	\node[draw=black!75, fill=black!6, line width=0.6pt, rounded corners=0.109cm, dashed, minimum width=2.520cm, minimum height=1.368cm, text width=2.412cm, align=center, inner sep=1pt, font=\diagfont] (14) at (8.460,-4.104) {Diffusion reference generation (removed: 3.9 s, 2.6 GB GPU)};
	\node[draw=black!75, fill=black!6, line width=0.6pt, rounded corners=0.081cm, minimum width=6.120cm, minimum height=1.008cm, text width=6.012cm, align=center, inner sep=1pt, font=\diagfont] (15) at (4.140,-9.864) {\textbf{EDICT hiding (baseline)} \\ IPAdapter + EDICT inversion + token mask \\ $\rightarrow$ $X_{\mathrm{stego}}$ (paper-style look)};
	\node[draw=black!75, fill=black!6, line width=0.6pt, rounded corners=0.081cm, minimum width=2.700cm, minimum height=1.008cm, text width=2.592cm, align=center, inner sep=1pt, font=\diagfont] (16) at (1.890,-12.024) {\textbf{Path A: semantic channel} \\ SwinJSCC + AWGN};
	\node[draw=black!75, fill=black!6, line width=0.6pt, rounded corners=0.270cm, minimum width=2.700cm, minimum height=0.540cm, text width=2.592cm, align=center, inner sep=1pt, font=\diagfont] (17) at (1.890,-13.590) {to receiver (Path A)};
	\node[draw=black!75, fill=black!6, line width=0.6pt, rounded corners=0.115cm, minimum width=3.060cm, minimum height=1.440cm, text width=2.952cm, align=center, inner sep=1pt, font=\diagfont] (7) at (10.530,-11.520) {\textbf{N6-N8: Shrink + lock the secret} \\ lossless WebP $\rightarrow$ encrypt with key + integrity seal (HMAC)};
	\draw[-{Stealth[length=1.8mm]}, line width=0.6pt] (4.950,-1.080) -- (4.950,-1.440) -- (3.420,-1.440) -- (3.420,-1.800);
	\draw[-{Stealth[length=1.8mm]}, line width=0.6pt] (4.950,-1.080) -- (4.950,-1.440) -- (7.830,-1.440) -- (7.830,-1.800);
	\draw[-{Stealth[length=1.8mm]}, line width=0.6pt] (3.420,-2.592) -- (3.420,-3.960);
	\draw[-{Stealth[length=1.8mm]}, line width=0.6pt] (6.480,-2.664) -- (6.480,-6.102) -- (5.760,-6.102);
	\draw[-{Stealth[length=1.8mm]}, line width=0.6pt] (3.420,-5.148) -- (3.420,-5.508);
	\draw[-{Stealth[length=1.8mm]}, line width=0.6pt] (3.420,-6.696) -- (3.420,-7.056);
	\draw[-{Stealth[length=1.8mm]}, line width=0.6pt] (3.420,-8.064) -- (3.420,-9.360);
	\draw[-{Stealth[length=1.8mm]}, line width=0.6pt] (1.890,-2.196) -- (0.450,-2.196) -- (0.450,-9.630) -- (1.080,-9.630);
	\draw[-{Stealth[length=1.8mm]}, line width=0.6pt] (3.600,-0.810) -- (0.180,-0.810) -- (0.180,-10.080) -- (1.080,-10.080);
	\draw[-{Stealth[length=1.8mm]}, line width=0.6pt] (10.170,-1.080) -- (10.170,-7.560) -- (5.760,-7.560);
	\draw[-{Stealth[length=1.8mm]}, line width=0.6pt] (10.530,-1.080) -- (10.530,-9.864) -- (7.200,-9.864);
	\draw[-{Stealth[length=1.8mm]}, line width=0.6pt] (10.620,-1.080) -- (11.160,-1.080) -- (11.160,-10.800);
	\draw[-{Stealth[length=1.8mm]}, line width=0.6pt] (4.950,-0.540) -- (4.950,-0.180) -- (12.420,-0.180) -- (12.420,-11.520) -- (12.060,-11.520);
	\draw[-{Stealth[length=1.8mm]}, line width=0.6pt] (10.530,-12.240) -- (10.530,-14.701) -- (8.280,-14.701);
	\draw[-{Stealth[length=1.8mm]}, line width=0.6pt] (6.300,-10.368) -- (6.300,-14.160);
	\draw[-{Stealth[length=1.8mm]}, line width=0.6pt] (6.300,-15.242) -- (6.300,-16.594);
	\draw[-{Stealth[length=1.8mm]}, line width=0.6pt] (1.890,-10.368) -- (1.890,-11.520);
	\draw[-{Stealth[length=1.8mm]}, line width=0.6pt] (1.890,-12.528) -- (1.890,-13.320);
	\node[fill=white, inner sep=1pt, font=\diaglabel, text=black!75] at (3.420,-3.564) {RACS reference acquisition};
	\node[anchor=east, fill=white, inner sep=1pt, font=\diaglabel, text=black!75] at (12.096,-17.550) {Exact-recovery layer (Path B)};
\end{tikzpicture}

	\caption{RACS sender (hiding).}
	\label{fig:racs-tx}
\end{figure}

\begin{figure}[H]
	\centering
	% Generated from racs_receiver.drawio by drawio2tikz.py; edit the .drawio file and rerun.
\begin{tikzpicture}
	\hyphenpenalty=10000\exhyphenpenalty=10000
	\draw[dashed, draw=black!70, rounded corners=2pt] (0.340,-2.210) rectangle (7.480,-7.990);
	\draw[dashed, draw=black!70, rounded corners=2pt] (8.160,-2.210) rectangle (12.750,-9.605);
	\node[draw=black!75, fill=black!6, line width=0.6pt, rounded corners=0.068cm, minimum width=2.550cm, minimum height=0.850cm, text width=2.448cm, align=center, inner sep=1pt, font=\diagfont] (dec) at (1.955,-3.536) {\textbf{SwinJSCC decoder}};
	\node[draw=black!75, fill=black!6, line width=0.6pt, rounded corners=0.117cm, minimum width=2.890cm, minimum height=1.462cm, text width=2.788cm, align=center, inner sep=1pt, font=\diagfont] (ret) at (5.695,-3.536) {\textbf{Same RACS retrieval} \\ CLIP search $\rightarrow$ rerank $\rightarrow$ token pick $\rightarrow$ identical $X_{\mathrm{ref}}$ (66/66)};
	\node[draw=black!75, fill=black!6, line width=0.6pt, rounded corners=0.076cm, minimum width=4.760cm, minimum height=0.952cm, text width=4.658cm, align=center, inner sep=1pt, font=\diagfont] (rec) at (3.740,-5.576) {\textbf{EDICT recovery + CodeFormer} \\ (baseline)};
	\node[draw=black!75, fill=black!6, line width=0.6pt, rounded corners=0.340cm, minimum width=4.760cm, minimum height=0.680cm, text width=4.658cm, align=center, inner sep=1pt, font=\diagfont] (oa) at (3.740,-7.140) {Recovered $\approx$ $X_{\mathrm{s}}$ (about 22.6 dB)};
	\node[draw=black!75, fill=black!6, line width=0.6pt, rounded corners=0.095cm, minimum width=3.910cm, minimum height=1.190cm, text width=3.808cm, align=center, inner sep=1pt, font=\diagfont] (find) at (10.455,-3.400) {\textbf{N12-N14: Find and read} \\ recompute busy-pixel spots from the image + key; read the $\pm$1 pattern};
	\node[draw=black!75, fill=black!6, line width=0.6pt, rounded corners=0.082cm, minimum width=3.230cm, minimum height=1.020cm, text width=3.128cm, align=center, inner sep=1pt, font=\diagfont] (seal) at (10.455,-5.185) {\textbf{N15: Check the seal} \\ HMAC with the key};
	\node[draw=black!75, fill=black!6, line width=0.6pt, rounded corners=0.082cm, minimum width=2.210cm, minimum height=1.020cm, text width=2.108cm, align=center, inner sep=1pt, font=\diagfont] (unl) at (9.520,-7.225) {\textbf{N16: Unlock} \\ decrypt with the key};
	\node[draw=black!75, fill=black!6, line width=0.6pt, rounded corners=0.095cm, dashed, minimum width=1.700cm, minimum height=1.190cm, text width=1.598cm, align=center, inner sep=1pt, font=\diagfont] (rej) at (11.730,-7.310) {Wrong key / modified image $\rightarrow$ rejected};
	\node[draw=black!75, fill=black!6, line width=0.6pt, rounded corners=0.340cm, minimum width=2.380cm, minimum height=0.680cm, text width=2.278cm, align=center, inner sep=1pt, font=\diagfont] (ob) at (9.520,-8.755) {\textbf{Recovered = $X_{\mathrm{s}}$ exactly}};
	\node[draw=black!75, fill=black!6, line width=0.6pt, rounded corners=0.340cm, minimum width=2.550cm, minimum height=0.680cm, text width=2.448cm, align=center, inner sep=1pt, font=\diagfont] (recv) at (1.955,-0.850) {Received signal (Path A)};
	\node[draw=black!75, fill=black!6, line width=0.6pt, rounded corners=0.340cm, minimum width=2.890cm, minimum height=0.680cm, text width=2.788cm, align=center, inner sep=1pt, font=\diagfont] (pub) at (5.695,-0.850) {$K_{\mathrm{pub}}$ + $X_{\mathrm{feat}}$ (public)};
	\node[draw=black!75, fill=black!6, line width=0.6pt, rounded corners=0.340cm, minimum width=2.040cm, minimum height=0.680cm, text width=1.938cm, align=center, inner sep=1pt, font=\diagfont] (tau) at (9.520,-0.850) {Private token $\tau$};
	\node[draw=black!75, fill=black!6, line width=0.6pt, rounded corners=0.340cm, minimum width=1.870cm, minimum height=0.680cm, text width=1.768cm, align=center, inner sep=1pt, font=\diagfont] (png) at (11.815,-0.850) {Stego PNG file (Path B)};
	\node[fill=white, inner sep=1pt, font=\diagfont] at (9.979,-6.018) {valid};
	\node[fill=white, inner sep=1pt, font=\diagfont] at (11.127,-6.018) {invalid};
	\draw[-{Stealth[length=1.8mm]}, line width=0.6pt] (1.955,-1.190) -- (1.955,-3.111);
	\draw[-{Stealth[length=1.8mm]}, line width=0.6pt] (5.695,-1.190) -- (5.695,-2.805);
	\draw[-{Stealth[length=1.8mm]}, line width=0.6pt] (8.500,-0.850) -- (7.820,-0.850) -- (7.820,-3.536) -- (7.140,-3.536);
	\draw[-{Stealth[length=1.8mm]}, line width=0.6pt] (1.955,-3.961) -- (1.955,-5.100);
	\draw[-{Stealth[length=1.8mm]}, line width=0.6pt] (5.695,-4.267) -- (5.695,-5.100);
	\draw[-{Stealth[length=1.8mm]}, line width=0.6pt] (3.740,-6.052) -- (3.740,-6.800);
	\draw[-{Stealth[length=1.8mm]}, line width=0.6pt] (9.520,-1.190) -- (9.520,-2.805);
	\draw[-{Stealth[length=1.8mm]}, line width=0.6pt] (11.815,-1.190) -- (11.815,-2.805);
	\draw[-{Stealth[length=1.8mm]}, line width=0.6pt] (10.455,-3.995) -- (10.455,-4.675);
	\draw[-{Stealth[length=1.8mm]}, line width=0.6pt] (10.455,-5.695) -- (10.455,-6.205) -- (9.520,-6.205) -- (9.520,-6.715);
	\draw[-{Stealth[length=1.8mm]}, line width=0.6pt] (10.455,-5.695) -- (10.455,-6.205) -- (11.730,-6.205) -- (11.730,-6.715);
	\draw[-{Stealth[length=1.8mm]}, line width=0.6pt] (9.520,-7.735) -- (9.520,-8.415);
	\node[anchor=west, fill=white, inner sep=1pt, font=\diaglabel, text=black!75] at (0.476,-7.871) {Path A: semantic channel (approximate)};
	\node[anchor=east, fill=white, inner sep=1pt, font=\diaglabel, text=black!75] at (12.614,-9.486) {Path B: exact recovery (no AI model)};
\end{tikzpicture}

	\caption{RACS receiver (recovery).}
	\label{fig:racs-rx}
\end{figure}

\paragraph{The reference is found, not painted.} The public key $K_{\mathrm{pub}}$ is turned into
a CLIP embedding and used to search a fixed collection of 1{,}000 images. The 50 closest images
are re-scored by how well they match in meaning (75\,\%) and in layout (25\,\%, comparing edge
maps with the public structure $X_{\mathrm{feat}}$), and the token picks one of the best 10. The
receiver repeats exactly the same steps from public data and the token and gets exactly the same
image (66 of 66 test images), in 0.05\,s and without a GPU.

\paragraph{The secret arrives exactly.} The stego image still travels through the semantic
channel (Path A), giving an approximate recovery as in the baseline. In addition, the original
secret is compressed losslessly, encrypted and sealed with a key derived from the token, and
hidden with tiny $\pm1$ changes in the most textured pixels of the stego image, which is sent as a
PNG file (Path B). The receiver recomputes the same pixel positions, checks the seal and decrypts.
A wrong key or an altered file is rejected instead of producing a misleading image.

\section{New-Block Reference}
\begin{table}[!htbp]
	\centering
	\caption{New RACS blocks (labels as in Figures~\ref{fig:racs-tx} and~\ref{fig:racs-rx}) and the reason for each design choice.}
	\label{tab:newblocks}
	\begin{tabularx}{\textwidth}{|c|l|X|}
		\hline
		\textbf{Block} & \textbf{What it does} & \textbf{Design choice and reason} \\ \hline
		N1 & CLIP text embedding of $K_{\mathrm{pub}}$ & Query from \emph{public} data only, so the receiver can recompute it without the secret. \\ \hline
		N2 & Exact search, top-50 pool & No approximate index, so sender and receiver get identical rankings. \\ \hline
		N3 & Edge-map structural score & CLIP finds the right content but not the right layout; $X_{\mathrm{feat}}$ is the only public structure both sides hold. \\ \hline
		N4 & 0.75 / 0.25 rerank, top 10 & Weights chosen on calibration images only; reranking a wider pool changes \emph{which} images the token can pick. \\ \hline
		N5 & SHA-256($\tau$) mod $K$ & One implementation at both ends; a 1-in-$K$ guess for an attacker. \\ \hline
		N6 & Lossless WebP compaction & 32\,\% smaller payload, which makes a low embedding rate affordable. \\ \hline
		N7 & PBKDF2 key derivation \cite{rfc8018} & Slows brute force of the short numeric token. \\ \hline
		N8, N16 & Stream encryption \cite{fips202} & Hidden bits look random. \\ \hline
		N8, N15 & HMAC-SHA256 seal \cite{rfc2104} & A wrong key or modified image is \emph{rejected}, never shown as a wrong image. \\ \hline
		N9 & Carrier upscale + dither & An upscaled image has an unnaturally clean LSB plane; camera-like noise makes it realistic. \\ \hline
		N10, N13 & Texture-adaptive $\pm1$ & Changes only busy regions; $\pm1$ instead of overwriting defeats the chi-square attack; the receiver recomputes positions with no side information. \\ \hline
		N11, N12 & PNG file transfer & Lossless; recompression would destroy the payload (and is detected). \\ \hline
	\end{tabularx}
\end{table}

\paragraph{One change at a time.} In the code, RACS is a subclass of the baseline that
overrides a single method: where the reference comes from. Every other stage is imported, not
copied, so any measured difference comes from the reference alone.

\section{Prototype Architecture}
\begin{figure}[H]
	\centering
	\begin{tikzpicture}[scale=0.85, transform shape]
		\node[src, text width=2.4cm] (s) at (0,0) {Sender phone\\(browser)};
		\node[blk, text width=3.0cm] (srv) at (5.2,0) {Laptop web server\\(runs the pipeline)};
		\node[src, text width=2.4cm] (r) at (10.4,0) {Receiver phone\\(browser)};
		\node[newblk, text width=3.4cm] (w) at (5.2,-2.0) {Worker process\\public key + RACS retrieval};
		\draw[arr] (s) -- node[above, font=\scriptsize]{secret photo} (srv);
		\draw[arr] (srv) -- node[above, font=\scriptsize]{stego PNG + key} (r);
		\draw[arr, <->] (srv) -- (w);
	\end{tikzpicture}
	\caption{Phone prototype: phones connect over Wi-Fi and the laptop does all the work.}
	\label{fig:app}
\end{figure}

The prototype offers three modes: \emph{hybrid} (default: diffusion stego look plus the exact
layer), \emph{lossless} (exact layer on a retrieved cover, no diffusion) and \emph{diffusion}
(the paper's approximate pipeline). The stego must be sent as a file, because messaging apps
recompress photos; the receiver detects this through the seal.

\section{Code Organisation}
\begin{table}[H]
	\centering
	\caption{Modules of the implementation.}
	\label{tab:modules}
	\begin{tabularx}{\textwidth}{|l|X|}
		\hline
		\textbf{Package} & \textbf{Responsibility} \\ \hline
		\texttt{common} & Paths, seeding, device selection, low-memory model loading \\ \hline
		\texttt{configs} & One configuration file read by both ends \\ \hline
		\texttt{semantic} & Public key and structure extraction \\ \hline
		\texttt{reference} & Baseline reference generation and ControlNet fine-tuning \\ \hline
		\texttt{retrieval} & RACS: corpus, search, structure cache, rerank, token pick \\ \hline
		\texttt{steganography} & EDICT, hiding, recovery, token mask, IPAdapter \\ \hline
		\texttt{communication} & AWGN channel and SwinJSCC codec \\ \hline
		\texttt{enhancement} & CodeFormer face restoration \\ \hline
		\texttt{pipelines} & The assembled baseline and RACS systems \\ \hline
		\texttt{evaluation} & Metrics, frozen test splits, reports \\ \hline
		\texttt{security} & Token tests, XuNet, chi-square and SPA detectors \\ \hline
		\texttt{app} & Phone web server and the exact-recovery layer \\ \hline
		\texttt{scripts} & Build and verification script for every phase \\ \hline
	\end{tabularx}
\end{table}
